\documentclass[11pt]{article}
\usepackage{mathblatt}
% gesetzt von werkzeuge/setzer.py (Sorte selbst) aus dem Lernweg TCQ
% und den Bankzeilen; Muster bau/proben/2026-10-09/hypotenuse-muster.
\usepackage{enumitem}
\usepackage{adjustbox,varwidth}
\newgeometry{a4paper,margin=18mm,top=15mm,bottom=20mm,footskip=9mm}
\pagestyle{fancy}\fancyhf{}\renewcommand{\headrulewidth}{0pt}
\fancyfoot[L]{\footnotesize\color{mbgrau}TCQ \textperiodcentered{} Lösungen \textperiodcentered{} Wachstum: linear oder exponentiell?}
\fancyfoot[R]{\footnotesize\color{mbgrau}\thepage}
\setlength{\parskip}{3pt}
\renewcommand{\kreuz}[1]{\mbox{$\square$\ #1}\quad}
\newcommand{\kreuzl}[1]{\par\hangindent1.4em\hangafter1\noindent$\square$\ #1}
\newcommand{\aufg}[3]{\par\noindent\makebox[0pt][r]{#3}\begin{minipage}[t]{\linewidth}#1\end{minipage}\par}
\newcommand{\aufgg}[4]{\par\noindent\makebox[0pt][r]{#4}\begin{minipage}[t]{0.6\linewidth}\vspace{0pt}#1\end{minipage}\hfill
  \begin{minipage}[t]{0.37\linewidth}\vspace{0pt}\centering #2\end{minipage}\par}
\newcommand{\nr}[1]{\makebox[7mm][l]{\textbf{#1}}}
\newcommand{\lsg}[3]{\par\noindent\begin{minipage}[t]{9mm}\textbf{#1}\end{minipage}%
  \begin{minipage}[t]{52mm}\raggedright\bfseries\boldmath #2\end{minipage}\hspace{3mm}%
  \begin{minipage}[t]{\dimexpr\linewidth-64mm\relax}\raggedright\small #3\end{minipage}\par\vspace{4pt}}
\newlength{\kr}
\makeatletter
\newcommand{\karomerk}[2]{\protected@write\@auxout{}{\string\karoseite{#1}{#2}{\thepage}}}
\newcommand{\karoseite}[3]{\expandafter\xdef\csname karo@#1@#2\endcsname{#3}}
\newcommand{\karohinweis}[3]{\karomerk{h}{#1}\ifcsname karo@k@#1\endcsname
  \ifnum\csname karo@k@#1\endcsname=0\csname karo@h@#1\endcsname\relax#2\else#3\fi
  \else#3\fi}
\makeatother
\begin{document}

\noindent{\Large\bfseries Lösungen}\hspace{1em}{\color{mbgrau}Wachstum: linear oder exponentiell?}\par\smallskip
\noindent{\small Links das Ergebnis, rechts der Weg in kurzen Schritten. Stimmt dein Ergebnis nicht, such den ersten Schritt, der anders ist.}\par
\par\Needspace{25mm}\vspace{6pt}\noindent\colorbox{black!12}{\makebox[7mm]{\bfseries\strut A}}\hspace{2mm}{\bfseries An der Tabelle erkennen}\par\vspace{4pt}
\lsg{1.}{linear: jede Minute $+30$ Liter}{Differenzen $30$, $30$, $30$: immer gleich, also linear; jede Minute $30$ Liter mehr.}
\lsg{2.}{exponentiell, Faktor $1{,}25$ ($+25\,\%$)}{Quotienten $400 : 320 = 1{,}25$, $500 : 400 = 1{,}25$, $625 : 500 = 1{,}25$ (Differenzen $80$, $100$, $125$): exponentiell, jedes Jahr $25\,\%$ mehr.}
\lsg{3.}{exponentiell, Faktor $0{,}8$ ($-20\,\%$)}{Quotienten $1\,600 : 2\,000 = 0{,}8$, $1\,280 : 1\,600 = 0{,}8$, $1\,024 : 1\,280 = 0{,}8$: exponentiell; Faktor $0{,}8$, also jedes Jahr $20\,\%$ weniger.}
\lsg{4.}{„um $10\,\%$“ und „Abnahme wird kleiner“}{Außerdem richtig: „Die Abnahme in mg wird jede Stunde kleiner.“ Quotienten immer $0{,}9$; Abnahmen $30$, $27$, $24{,}3$ mg.}
\par\Needspace{25mm}\vspace{6pt}\noindent\colorbox{black!12}{\makebox[7mm]{\bfseries\strut B}}\hspace{2mm}{\bfseries Am Text erkennen}\par\vspace{4pt}
\lsg{1.}{a) L ab \ b) E zu \ c) E ab \ d) L zu}{a) L, ab; b) E, zu; c) E, ab; d) L, zu}
\lsg{2.}{Nein; Verlust $3\,600\,€$, dann $3\,060\,€$}{Nein; $24\,000 \cdot 0{,}85 = 20\,400$, $20\,400 \cdot 0{,}85 = 17\,340$; Verlust $3\,600\,€$, dann nur $3\,060\,€$.}
\lsg{3.}{Nein; $900$ Nutzer, nicht $800$}{Nein; $400 \cdot 1{,}5 = 600$, $600 \cdot 1{,}5 = 900$; doppelt wären $800$.}
\lsg{4.}{Woche 4: A; B ab Woche $6$}{Woche 4: A $= 64\,€$, B $\approx 60{,}84\,€$, also A. Ab Woche 6 zahlt B mehr: $80{,}45\,€ > 80\,€$. (A linear: $40$, $48$, $56$, $64$, $72$, $80$. B exponentiell: $40$, $46$, $52{,}90$, $60{,}84$, $69{,}96$, $80{,}45$.)}
\par\Needspace{25mm}\vspace{6pt}\noindent\colorbox{black!12}{\makebox[7mm]{\bfseries\strut C}}\hspace{2mm}{\bfseries Den passenden Graphen finden}\par\vspace{4pt}
\lsg{1.}{$f$ Gerade, steigt, $100$; $g$ Kurve, fällt, $400$; $h$ Kurve, steigt, $50$}{Startwerte $100$, $400$, $50$. $f$: Gerade, steigt; $g$: Kurve, fällt; $h$: Kurve, steigt.}
\lsg{2.}{$g$}{jede Minute derselbe Betrag, also eine Gerade, und sie beginnt bei $150$.}
\lsg{3.}{$g$}{beginnt bei $400$ und wird immer steiler. $f$ nicht: Gerade, also jedes Jahr derselbe Betrag. $h$ nicht: beginnt bei $0$, nicht bei $400$.}
\par\Needspace{25mm}\vspace{6pt}\noindent\colorbox{black!12}{\makebox[7mm]{\bfseries\strut D}}\hspace{2mm}{\bfseries Wertepaare als Punkte eintragen}\par\vspace{4pt}
\lsg{1.}{Punkte $(0|10)$, $(1|20)$, $(2|40)$, $(3|80)$}{\par\smallskip \adjustbox{max width=\linewidth,max height=4cm}{\begin{varwidth}{50cm}\begin{ksys}[xmin=0,xmax=5,ymin=0,ymax=80,xstep=1,ystep=10,karo=0.35,xlabel=Monat,ylabel=Mitglieder] \punkt{0}{10}{} \punkt{1}{20}{} \punkt{2}{40}{} \punkt{3}{80}{} \end{ksys}\end{varwidth}}}
\lsg{2.}{Punkte $(0|40)$ bis $(4|154)$, Kurve}{Punkte $(0|40)$, $(1|56)$, $(2|78)$, $(3|110)$, $(4|154)$; glatte Kurve, immer steiler.\par\smallskip \adjustbox{max width=\linewidth,max height=4cm}{\begin{varwidth}{50cm}\begin{ksys}[xmin=0,xmax=5,ymin=0,ymax=160,xstep=1,ystep=20,karo=0.35,xlabel=Stunde,ylabel=Zellen] \punkt{0}{40}{} \punkt{1}{56}{} \punkt{2}{78}{} \punkt{3}{110}{} \punkt{4}{154}{} \funktionab{40*1.4^\x}{}{0}{4} \end{ksys}\end{varwidth}}}
\lsg{3.}{exponentiell: Punkte auf einer immer steileren Kurve}{Exponentiell; $y$: 1 Kästchen = 5 Tausend; Punkte $(0|8)$, $(1|12)$, $(2|18)$, $(3|27)$, $(4|40{,}5)$ liegen auf einer Kurve, die immer steiler wird (Quotient immer $1{,}5$).\par\smallskip \adjustbox{max width=\linewidth,max height=4cm}{\begin{varwidth}{50cm}\begin{ksys}[xmin=0,xmax=5,ymin=0,ymax=45,xstep=1,ystep=5,karo=0.3,xlabel=Woche,ylabel=Tausend] \punkt{0}{8}{} \punkt{1}{12}{} \punkt{2}{18}{} \punkt{3}{27}{} \punkt{4}{40.5}{} \funktionab{8*1.5^\x}{}{0}{4} \end{ksys}\end{varwidth}}}
\par\Needspace{25mm}\vspace{6pt}\noindent\colorbox{black!12}{\makebox[7mm]{\bfseries\strut T}}\hspace{2mm}{\bfseries Probetest}\par\vspace{4pt}
\lsg{1.}{exponentiell, Faktor $1{,}04$}{Exponentiell; jeder Wert ist das $1{,}04$-Fache des vorigen ($52\,000 : 50\,000 = 1{,}04$, $54\,080 : 52\,000 = 1{,}04$); die Zuwächse ($2\,000$, dann $2\,080$) werden größer.}
\lsg{2.}{$g$}{beginnt bei $800$, fällt und wird flacher. $f$ ist eine Gerade, $h$ beginnt bei $600$.}
\lsg{3.}{Punkte $(0|25)$ bis $(4|164)$, Kurve}{Punkte $(0|25)$, $(1|40)$, $(2|64)$, $(3|102)$, $(4|164)$; glatte Kurve.\par\smallskip \adjustbox{max width=\linewidth,max height=4cm}{\begin{varwidth}{50cm}\begin{ksys}[xmin=0,xmax=5,ymin=0,ymax=180,xstep=1,ystep=20,karo=0.3,xlabel=Stunde,ylabel=Bakterien] \punkt{0}{25}{} \punkt{1}{40}{} \punkt{2}{64}{} \punkt{3}{102}{} \punkt{4}{164}{} \funktionab{25*1.6^\x}{}{0}{4} \end{ksys}\end{varwidth}}}
\lsg{4.}{Nein; nach $6$ Wochen bedeckt}{Verdoppeln ist exponentiell: $3$, $6$, $12$, $24$, $48$, $96$, $192$ m². Nach $5$ Wochen $96$ m², nach $6$ Wochen wären es $192$ m² $> 120$ m². Danach ist kein Platz mehr.}
\end{document}
